\subsection{Debugging}
Large Language Models (LLMs) are increasingly being integrated into Automated Program Repair (APR) systems due to their strong capabilities in code understanding and generation \cite{llm_repair_1}. Recent research has shown that models such as Codex, CodeT5, and GPT can effectively identify and repair software bugs by leveraging contextual and semantic understanding of source code \cite{llm_repair_3,llm_repair_4}. Unlike traditional rule-based APR techniques, LLM-based approaches can generate flexible and context-aware patches for a variety of programming errors \cite{llm_repair_6}.

In a structured literature review done by \citet{zubair2025llmprogramrepair}, the authors mention that LLMs hold promise for automating program repair tasks with a combination of key characteristics. The most mentioned capabilities are code generation and contextual understanding. These capabilities allow LLMs to create and modify the code while considering its structure and meaning. Fine-tuning and multilingual support are also emerging research areas, indicating the potential of LLMs to specialize in specific repair tasks and handle different programming languages.

\subsubsection{What specific programming errors or defects are targeted in program repair tasks using LLMs?}
\citet{zubair2025llmprogramrepair} mentions that semantic and syntax errors are the most focused errors for program repair. There is a clear focus on software bugs when compared to security vulnerabilities. This indicates a stronger emphasis on software functionality repair rather than security-specific issues. The focus on single-hunk repair\footnote{Single-hunk bugs require changes in one contiguous code location \cite{zubair2025llmprogramrepair}} suggests that LLMs models currently are more focused on simpler code modifications. More investigation needs to be done on multi-hunk bugs\footnote{Multi-hunk bugs require fixes spanning multiple, non-contiguous locations, often in different files \cite{zubair2025llmprogramrepair}.}.

\subsubsection{Debugging challenges of LLMs} Although Large Language Models (LLMs) have shown promising capabilities in automated debugging, there exist several important limitations in their effectiveness. A study by \citet{tian2024debugbench} found that LLMs perform well on syntax and reference errors but struggle significantly with logical and multi-error bugs that require deeper semantic reasoning and understanding of execution flow. Runtime feedback, such as traceback messages, can improve fixes for simple errors but may negatively affect logical debugging by misleading the model toward superficial corrections. Additionally, even advanced models like GPT-4 \cite{openai_gpt4} still perform below experienced human programmers in complex debugging scenarios, indicating that current LLMs lack robust reasoning and reliability for real-world software debugging tasks.